
\section{Implementation for Image Generation}
\label{sec:impl}

We describe our implementation for image generation on ImageNet \cite{deng2009imagenet} at resolution 256$\times$256. Full implementation details are provided in Appendix \ref{sec:appendix_impl}.

\paragraph{Tokenizer.} By default, we perform generation in latent space \cite{rombach2022high}. We adopt the standard SD-VAE tokenizer, which produces a 32$\times$32$\times$4 latent space in which generation is performed.

\paragraph{Architecture.} Our generator ($f_\theta$) has a DiT-like \cite{peebles2023scalable} architecture.
Its input is 32$\times$32$\times$4-dim Gaussian noise $\e$, and its output is the generated latent $\x$ of the same dimension.
We use a patch size of 2, \ie, like DiT/2.
Our model uses adaLN-zero \cite{peebles2023scalable} for processing class-conditioning or other extra conditioning.


\paragraph{CFG conditioning.} We follow \cite{geng2025improved} and adopt CFG-conditioning. At training time, a CFG scale $\alpha$ (Eq.~(\ref{eq:cfg2})) is randomly sampled. Negative samples are prepared based on $\alpha$ (Eq.~(\ref{eq:cfg1})), and the network is conditioned on this value. At inference time, $\alpha$ can be freely specified and varied without retraining. Details are in \ref{sec:appendix_cfg}.


\paragraph{Batching.} The pseudo-code in Alg.~\ref{alg:training_code} describes a batch of $N=N_\text{neg}$ generated samples. In practice, when class labels are involved, we sample a batch of $N_\text{c}$ class labels. For each label, we perform Alg.~\ref{alg:training_code} \textit{independently}. Accordingly, the \textit{effective} batch size is $B=N_\text{c}{\times}N$, which consists of $N_\text{c}{\times}N$ negatives and $N_\text{c}{\times}N_\text{pos}$ positives.

We define a ``training epoch'' based on the number of generated samples $\x$.
In particular, each iteration generates $B$ samples, and one epoch corresponds to $N_\text{data}/B$ iterations for a dataset of size $N_\text{data}$.

\paragraph{Feature Extractor.} Our model is trained with drifting loss in a feature space (Sec.~\ref{subsec:feature_space}). 
The feature extractor $\phi$ is an image encoder.
We mainly consider a ResNet-style \cite{he2016deep} encoder, pre-trained by self-supervised learning, \eg, MoCo \cite{he2020momentum} and SimCLR \cite{chen2020simple}.
When these pre-trained models operate in pixel space, we apply the VAE decoder to map our generator’s latent-space output back to pixel space for feature extraction. Gradients are backpropagated through the feature encoder and VAE decoder.
We also study an MAE \cite{he2022masked} pre-trained in latent space (detailed in \ref{sec:appendix_latent_mae}).

For all ResNet-style models, features are extracted from multiple stages (\ie, multi-scale feature maps). The drifting loss in (\ref{eq:feature_space}) is computed at each scale and then combined. We elaborate on the details in \ref{sec:appendix_feature_and_loss_normalization}.


\paragraph{Pixel-space Generation.} 
While our experiments primarily focus on latent-space generation, our models support pixel-space generation. In this case, $\e$ and $\x$ are both 256$\times$256$\times$3. We use a patch size of 16 (\ie, DiT/16). The feature extractor $\phi$ is directly on the pixel space.


